%! Introducing Statistics: Why Be Normal? — Unit 6, Lesson 6.1 — Guided Notes (Answer Key)
\documentclass[10pt]{article}
\usepackage{apstats-article}
\usepackage{apstats-key}   % pulls in -boxes; do NOT also load -boxes
\newcommand{\TallMath}[1]{$\displaystyle #1\rule[-1.4em]{0pt}{3.2em}$}

\begin{document}

\pageheader{Unit 6, Lesson 6.1}{Guided Notes --- Answer Key}
\namedateperiod

\vspace{0.15in}

\begin{objectivebox}
By the end of this lesson, I will be able to:
\begin{enumerate}
  \item \ansline{identify whether a research question calls for a confidence interval or a significance test}
  \item \ansline{verify the conditions needed to use a normal model for inference on a population proportion}
\end{enumerate}
\end{objectivebox}

\vspace{0.1in}

\begin{vocabbox}
\termblanklong{Statistical inference} \ansline{using sample data to draw conclusions about an unknown population parameter}
\termblanklong{Confidence interval} \ansline{a range of plausible values for an unknown parameter, constructed from sample data}
\termblanklong{Significance test (hypothesis test)} \ansline{a procedure for deciding whether sample data provide sufficient evidence against a specific claim about a parameter}
\termblanklong{Null hypothesis ($H_0$)} \ansline{the default claim about the parameter, assumed true unless evidence suggests otherwise}
\termblanklong{Alternative hypothesis ($H_a$)} \ansline{the claim we seek statistical evidence for; it contradicts $H_0$}
\end{vocabbox}

\vspace{0.1in}

\begin{notesbox}{The Big Shift: From Unit 5 to Unit 6}
\textbf{Unit 5 question:} Given $p = 0.60$ and $n = 100$, what is the probability that
$\hat p > 0.65$?

\textbf{Unit 6 question:} Given $\hat p = 0.65$ and $n = 100$, what can we say about
the \emph{unknown} population proportion $p$?

The key: since the sampling distribution of $\hat p$ is approximately
\ans{normal} when conditions are met, we can \ansline{work backwards from $\hat p$ to estimate or test claims about the unknown parameter $p$}.

\end{notesbox}

\vspace{0.1in}

\begin{notesbox}{Two Types of Inference}
\renewcommand{\arraystretch}{1.8}
\begin{tabular}{@{} >{\bfseries}p{0.20\linewidth} p{0.36\linewidth} p{0.36\linewidth} @{}}
  & \textbf{Confidence Interval} & \textbf{Significance Test} \\
  \hline
  Goal    & \ansline{Estimate the value of an unknown parameter} & \ansline{Decide if data provide evidence against a specific parameter value} \\
  Question asked & ``Where is \ans{$p$}?'' & ``Is \ans{$p_0$} consistent with my data?'' \\
  Answer gives & A range of \ans{plausible} values & A decision: reject or \ans{fail to reject} $H_0$ \\
\end{tabular}
\end{notesbox}

\vspace{0.1in}

\begin{notesbox}{Conditions for Using a Normal Model}
Before using any inference procedure in Unit 6, verify:

\begin{enumerate}
  \item \textbf{Random:} The data come from \ansline{a random sample or randomized experiment}.
  \item \textbf{Large Counts:} Both $n\hat{p} \ge$ \ans{10} and $n(1-\hat{p}) \ge$ \ans{10}.
  \item \textbf{Independence (10\% rule):} If sampling without replacement, $n \le$ \ans{10\% of $N$}.
\end{enumerate}

When all three conditions are met, the sampling distribution of $\hat p$ is approximately
\ans{normal} with mean \ans{$p$} and standard deviation \ans{$\sqrt{p(1-p)/n}$}.
\end{notesbox}

\vspace{0.1in}

\begin{practicebox}
\textbf{Quick Check.} Classify each research question as calling for a \emph{confidence interval}
(CI) or a \emph{significance test} (Test).

\begin{enumerate}
  \item ``What proportion of U.S. adults exercise daily?'' \hfill \ans{CI}
  \item ``Is there evidence that more than half of students prefer online homework?'' \hfill \ans{Test}
  \item ``We want to estimate, within 3 percentage points, the rate of vaccine hesitancy.''
        \hfill \ans{CI}
\end{enumerate}
\end{practicebox}

\end{document}
